%% =====================================================================
%%  B4-T-10-02 -- what B4 contributes to "The Reactance Trigger"
%%  -------------------------------------------------------------------
%%  LAYER A: APPEND-ONLY. Written once, then left alone. Material found
%%  only here sits in the `unique` slot and is protected by rule R10.
%% =====================================================================
%% @kind: source
%% @id: B4-T-10-02
%% @book: B4
%% @topic: T-10-02
%% @locator: ch. 7, pp. 205--210
%% @status: distilled
%% @unique: yes
%% @integrated: yes
%% @updated: 2026-09-19

\dossier{B4}{T-10-02}{\bookshort{B4}}{ch. 7, pp. 205--210}

\begin{distilled}
  \item Restricting a freedom increases the desire for it and for whatever it provided access to.
  \item The reaction defends the range of options rather than any single option, which is why the escalation exceeds the object's value.
  \item Withdrawal of an established freedom is far more powerful than refusal of a new one.
  \item Competition intensifies the effect by combining loss of an option with the presence of a rival.
\end{distilled}

\begin{unique}
  \item Reactance as the mechanism underneath scarcity, which explains why the value increase is displaced onto the object and survives the restriction.
  \item The asymmetry between withdrawal and refusal, and the resulting practical rule that granting-then-removing outperforms never offering.
\end{unique}
